\section{Background}
\label{sec:background}

\subsection{Microscaling formats}
\label{sec:bg-mx}

An MX block comprises $k$ elements sharing a single E8M0 scale
$2^{S-127}$, $S \in [0,254]$. The OCP specification fixes $k=32$; the
architectural granularity in this design is the 16-element half-block carried
by a 64-bit dual-read operand pair. Elements are E2M1 for \fpfour{} and E4M3 or
E5M2 for \fpeight. For two blocks $A$, $B$ with scales $s_a=2^{S_a-127}$,
$s_b=2^{S_b-127}$ accumulating into an FP32 value, the scaled dot product is
%
\begin{equation}
  \label{eq:base}
  \mathit{acc}' = s_a s_b \sum_i a_i b_i + \mathit{acc}.
\end{equation}
%
Both scales are powers of two, so applying them is an exponent adjustment
rather than a multiplication. Rewriting~\eqref{eq:base} to expose this,
%
\begin{equation}
  \label{eq:factor}
  \mathit{acc}' = s_a s_b\!\left( \sum_i a_i b_i
                  + \frac{\mathit{acc}}{s_a s_b} \right),
\end{equation}
%
places the products at a fixed fixed-point position and admits the accumulator
by a shift of $E_{\mathit{acc}} - \mathrm{bias} - (S_a{+}S_b{-}254)$ bit
positions, where $E_{\mathit{acc}}$ is the accumulator's biased exponent. No
multiplier is required on the datapath; the block scale reappears only in the
final exponent. Equation~\eqref{eq:factor} is the identity every engine
realises.

\subsection{Coprocessor interface}
\label{sec:bg-xif}

The host core, \core, is designed within OpenHW Group's CORE-V family to
extend the ISA through an external coprocessor over \xif{} rather than by
embedding custom functional units in its own pipeline, which keeps the
coprocessor's cost isolated and measurable rather than conflated with changes
to the host itself. The coprocessor attaches over \xif~\cite{openhw_cvxif} and
uses only the standardised issue, commit, and result channels, with no private
path into the host pipeline; any XIF compliant core supporting three source
reads and dual-read can host it. The host in this work, \core, does not
implement \xif's optional dual-read, which was added to its register-file read
path; the modification is independent of the coprocessor, which uses only the
standard interface.

\subsection{Motivation for \fpfourres}
\label{sec:bg-motiv}

\fpfour{} halves weight footprint versus \fpeight{} but at a real cost in accuracy~\cite{m2xfp2026}; uniform \fpeight{} recovers that accuracy but surrenders the memory-bandwidth advantage and widens the accumulation datapath. Asymmetric precision offers a middle path. On GPT-2/WikiText-2, naive \fpfour{} collapses perplexity from 30.3 to 107.8; adding a second \fpfour{} block to activations alone (\fpfourres, averaging 6.375\,b/element) recovers \SI{93}{\percent} of that loss, against \SI{36}{\percent} for the mirrored weight-side allocation. This activation-side advantage is not universal, however: \S\ref{sec:analysis} shows where it holds, where it inverts, and why.